\documentclass[11pt]{article}

\usepackage[a4paper,margin=1in]{geometry}
\usepackage{amsmath,amssymb,bm}

\title{Analytic Born--Oppenheimer Hamiltonian Derivatives}
\author{pyqed notes}
\date{}

\begin{document}
\maketitle

\section*{Born--Oppenheimer Hamiltonian}

The electronic Born--Oppenheimer Hamiltonian is
\begin{equation}
\hat H_{\mathrm{BO}}(\bm r; \bm R)
=
\hat T_e + \hat V_{ee} + \hat V_{en}(\bm R) + V_{nn}(\bm R),
\end{equation}
with
\begin{equation}
\hat V_{en}(\bm R)= -\sum_{iA}\frac{Z_A}{|\bm r_i-\bm R_A|}
\end{equation}
and
\begin{equation}
V_{nn}(\bm R)=\sum_{A<B}\frac{Z_A Z_B}{|\bm R_A-\bm R_B|}.
\end{equation}

Since only $\hat V_{en}$ and $V_{nn}$ depend explicitly on nuclear coordinates,
\begin{equation}
\hat F_k
=
\frac{\partial \hat H_{\mathrm{BO}}}{\partial Q_k}
=
\frac{\partial \hat V_{en}}{\partial Q_k}
+
\frac{\partial V_{nn}}{\partial Q_k}\hat I,
\end{equation}
and
\begin{equation}
\hat G_{kl}
=
\frac{\partial^2 \hat H_{\mathrm{BO}}}{\partial Q_k \partial Q_l}
=
\frac{\partial^2 \hat V_{en}}{\partial Q_k \partial Q_l}
+
\frac{\partial^2 V_{nn}}{\partial Q_k \partial Q_l}\hat I.
\end{equation}

\section*{Cartesian Nuclear Derivatives}

Using Cartesian nuclear coordinates $R_{A\alpha}$, with atom index $A$ and
Cartesian component $\alpha \in \{x,y,z\}$,
\begin{equation}
\hat F_{A\alpha}
=
\frac{\partial \hat H_{\mathrm{BO}}}{\partial R_{A\alpha}},
\qquad
\hat G_{A\alpha,B\beta}
=
\frac{\partial^2 \hat H_{\mathrm{BO}}}
{\partial R_{A\alpha}\partial R_{B\beta}}.
\end{equation}

Let
\begin{equation}
\bm x = \bm r - \bm R_A,
\qquad
x_\alpha = r_\alpha - R_{A\alpha},
\qquad
\rho = |\bm x|.
\end{equation}
Then
\begin{equation}
\frac{1}{|\bm r-\bm R_A|} = \frac{1}{\rho},
\qquad
\frac{\partial}{\partial R_{A\alpha}} = -\frac{\partial}{\partial x_\alpha}.
\end{equation}

Therefore, the first derivative is
\begin{equation}
\frac{\partial}{\partial R_{A\alpha}} \frac{1}{\rho}
=
\frac{x_\alpha}{\rho^3}
=
\frac{r_\alpha - R_{A\alpha}}{|\bm r-\bm R_A|^3}.
\end{equation}

For the electron--nuclear operator,
\begin{equation}
\frac{\partial \hat V_{en}}{\partial R_{A\alpha}}
=
- Z_A \sum_i \frac{r_{i\alpha}-R_{A\alpha}}{|\bm r_i-\bm R_A|^3}.
\end{equation}

The second derivative is
\begin{equation}
\frac{\partial^2}{\partial R_{A\alpha}\partial R_{A\beta}} \frac{1}{\rho}
=
-\frac{\delta_{\alpha\beta}}{\rho^3}
+
\frac{3 x_\alpha x_\beta}{\rho^5},
\end{equation}
or equivalently,
\begin{equation}
\frac{\partial^2}{\partial R_{A\alpha}\partial R_{A\beta}}
\frac{1}{|\bm r-\bm R_A|}
=
-\frac{\delta_{\alpha\beta}}{|\bm r-\bm R_A|^3}
+
\frac{3(r_\alpha-R_{A\alpha})(r_\beta-R_{A\beta})}{|\bm r-\bm R_A|^5}.
\end{equation}

Thus,
\begin{equation}
\frac{\partial^2 \hat V_{en}}
{\partial R_{A\alpha}\partial R_{A\beta}}
=
- Z_A \sum_i
\left[
-\frac{\delta_{\alpha\beta}}{|\bm r_i-\bm R_A|^3}
+
\frac{3(r_{i\alpha}-R_{A\alpha})(r_{i\beta}-R_{A\beta})}
{|\bm r_i-\bm R_A|^5}
\right].
\end{equation}

For different nuclei,
\begin{equation}
\frac{\partial^2 \hat V_{en}}
{\partial R_{A\alpha}\partial R_{B\beta}} = 0,
\qquad A \ne B.
\end{equation}

\section*{AO Matrix Elements}

In an AO basis $\{\mu,\nu\}$, define
\begin{equation}
h^{(1)}_{A\alpha,\mu\nu}
=
\left\langle \mu \left|
\frac{\partial \hat V_{en}}{\partial R_{A\alpha}}
\right| \nu \right\rangle
\end{equation}
and
\begin{equation}
h^{(2)}_{A\alpha,B\beta,\mu\nu}
=
\left\langle \mu \left|
\frac{\partial^2 \hat V_{en}}
{\partial R_{A\alpha}\partial R_{B\beta}}
\right| \nu \right\rangle.
\end{equation}

\section*{State-Basis Matrix Elements}

Let the spin-traced one-particle transition density matrix be
\begin{equation}
\gamma^{\beta\alpha'}_{\nu\mu}
=
\langle \Phi_\beta | a_\mu^\dagger a_\nu | \Phi_{\alpha'} \rangle.
\end{equation}
Then
\begin{equation}
F_{A\alpha}^{\beta\alpha'}
=
\langle \Phi_\beta | \hat F_{A\alpha} | \Phi_{\alpha'} \rangle
=
\sum_{\mu\nu}
\gamma^{\beta\alpha'}_{\nu\mu}\,
h^{(1)}_{A\alpha,\mu\nu}
+
\delta_{\beta\alpha'}
\frac{\partial V_{nn}}{\partial R_{A\alpha}},
\end{equation}
and
\begin{equation}
G_{A\alpha,B\beta}^{\beta\alpha'}
=
\langle \Phi_\beta | \hat G_{A\alpha,B\beta} | \Phi_{\alpha'} \rangle
=
\sum_{\mu\nu}
\gamma^{\beta\alpha'}_{\nu\mu}\,
h^{(2)}_{A\alpha,B\beta,\mu\nu}
+
\delta_{\beta\alpha'}
\frac{\partial^2 V_{nn}}
{\partial R_{A\alpha}\partial R_{B\beta}}.
\end{equation}

\section*{Projection to Normal Coordinates}

For a linear coordinate transform
\begin{equation}
Q_k = \sum_{A\alpha} L_{k,A\alpha} R_{A\alpha},
\end{equation}
the projected first- and second-order derivative matrices are
\begin{equation}
F_k^{\beta\alpha'}
=
\sum_{A\alpha} L_{k,A\alpha}\,
F_{A\alpha}^{\beta\alpha'},
\end{equation}
and
\begin{equation}
G_{kl}^{\beta\alpha'}
=
\sum_{A\alpha,B\beta}
L_{k,A\alpha}L_{l,B\beta}\,
G_{A\alpha,B\beta}^{\beta\alpha'}.
\end{equation}

\end{document}
